\documentclass[letterpaper,11pt]{article}

\usepackage{latexsym}
\usepackage[empty]{fullpage}
\usepackage{titlesec}
\usepackage{marvosym}
\usepackage[usenames,dvipsnames]{color}
\usepackage{verbatim}
\usepackage{enumitem}
\usepackage[hidelinks]{hyperref}
\usepackage{fancyhdr}
\usepackage[english]{babel}
\usepackage{tabularx}
\usepackage{fontawesome5}
\ifdefined\pdfgentounicode
  \input{glyphtounicode}
\fi

\pagestyle{fancy}
\fancyhf{}
\fancyfoot{}
\renewcommand{\headrulewidth}{0pt}
\renewcommand{\footrulewidth}{0pt}
\setlength{\footskip}{4.1pt}

\addtolength{\oddsidemargin}{-0.55in}
\addtolength{\evensidemargin}{-0.5in}
\addtolength{\textwidth}{1.15in}
\addtolength{\topmargin}{-0.5in}
\addtolength{\textheight}{1.35in}

\urlstyle{same}
\raggedbottom
\raggedright
\setlength{\tabcolsep}{0in}

\titleformat{\section}{
  \vspace{-4pt}\scshape\raggedright\large\bfseries
}{}{0em}{}[\color{black}\titlerule \vspace{-4pt}]

\ifdefined\pdfgentounicode
  \pdfgentounicode=1
\fi

\newcommand{\resumeItem}[1]{
  \item\small{{#1 \vspace{-2pt}}}
}

\newcommand{\resumeSubheading}[4]{
  \vspace{-1pt}\item
    \begin{tabular*}{1.0\textwidth}[t]{l@{\extracolsep{\fill}}r}
      \textbf{#1} & \textbf{\small #2} \\
      \textit{\small#3} & \textit{\small #4} \\
    \end{tabular*}\vspace{-6pt}
}

\newcommand{\resumeProjectHeading}[2]{
    \item
    \begin{tabular*}{1.0\textwidth}{l@{\extracolsep{\fill}}r}
      \small#1 & \textbf{\small #2}\\
    \end{tabular*}\vspace{-6pt}
}

\renewcommand\labelitemi{$\vcenter{\hbox{\tiny$\bullet$}}$}
\renewcommand\labelitemii{$\vcenter{\hbox{\tiny$\bullet$}}$}

\newcommand{\resumeSubHeadingListStart}{\begin{itemize}[leftmargin=0.0in, label={}]}
\newcommand{\resumeSubHeadingListEnd}{\end{itemize}}
\newcommand{\resumeItemListStart}{\begin{itemize}}
\newcommand{\resumeItemListEnd}{\end{itemize}\vspace{-4pt}}

\begin{document}

\begin{center}
    {\Huge \scshape \textbf{Martin Wu}} \\ \vspace{2pt}
    Vancouver, BC $|$
    \href{mailto:martinwu500@gmail.com}{\underline{martinwu500@gmail.com}} $|$
    \href{https://martinwu.dev}{\underline{martinwu.dev}} $|$
    \href{https://www.linkedin.com/in/martinwuu/}{\underline{linkedin.com/in/martinwuu}} $|$
    \href{https://github.com/martinw500}{\underline{github.com/martinw500}}
\end{center}

\section{Education}
  \resumeSubHeadingListStart
    \resumeSubheading
      {University of British Columbia}{Expected May 2029}
      {Bachelor of Applied Science in Computer Engineering}{Internship availability: Jan or May 2027, 4, 8 or 12 months}
      \resumeItemListStart
        \resumeItem{Average: A- $|$ Coursework: Computer Systems (SystemVerilog, FPGA, RISC-V), Software Construction (Java)}
      \resumeItemListEnd
  \resumeSubHeadingListEnd

\section{Experience}
  \resumeSubHeadingListStart

    \resumeSubheading
      {Embedded Systems Developer}{Sep 2025 -- Present}
      {UBC Solar}{Vancouver, BC}
      \resumeItemListStart
        \resumeItem{Brought up the Motor Drive Interface firmware on STM32: pin setup, a CAN driver with a hardware filter for motor commands, and ADC readings of board voltage and temperature}
        \resumeItem{Made motor commands expire after 100\,ms of CAN silence, so a stale command cannot keep driving the motor; confirmed on hardware by cutting the command stream}
        \resumeItem{Proposed PT1000 RTD motor-temperature sensing with a MAX31865 and wrote its firmware from the first prototype to the SPI driver; verified with heat-gun tests and integration-tested with the car's electrical system}
        \resumeItem{Built RTD fault handling: a leaky-counter debounce latches after 5 bad reads, fault and SPI-error bits go out on CAN, and a config read-back catches a missing chip that would otherwise read $-$259\,°C}
        \resumeItem{Wrote external-lights firmware as FreeRTOS tasks (turn signals, hazards, brake lights) and steering-wheel code that shows CAN speed data in km/h or mph on the hex display}
        \resumeItem{Added the new CAN signals to the telemetry DBC, the Grafana dashboards and the pit-crew test, so the crew sees motor-board data and sensor faults live}
      \resumeItemListEnd

    \resumeSubheading
      {Co-Founder \& Software Engineer}{Oct 2025 -- Present}
      {Stealth AI Infrastructure Startup}{Vancouver, BC}
      \resumeItemListStart
        \resumeItem{Wrote a Go test harness whose in-memory adapters drive reads, idempotent replays, identity denials and adapter panics through the team's 11-step dispatcher without Docker}
        \resumeItem{Fixed the tooling and rewrote the demo script for the live product demo, where AI agents query Postgres through access policies, human approvals and an audit trail; it is used in investor meetings and accelerator applications}
        \resumeItem{Ran the technical review for the first enterprise deal, an on-prem AI cluster now at the LOI stage: showed the model would not fit unquantized, flagged support and purchasing risks, and drafted the LOI}
      \resumeItemListEnd

  \resumeSubHeadingListEnd

\section{Projects}
    \resumeSubHeadingListStart
      \resumeProjectHeading
          {\textbf{Songless} $|$ \emph{React, TypeScript, Vite, Cloudflare R2} $|$ \href{https://songless-zeta.vercel.app}{Live}}{2026}
          \resumeItemListStart
            \resumeItem{Built a music-guessing web game with \textbf{100k+ visits} and 804 songs; 30-second clues stream from R2, and the full track loads only on the results screen}
            \resumeItem{Wrote the upload pipeline: trims leading silence, encodes 128\,kbps MP3s, and refuses any batch that would push the bucket past an 8.5\,GB ceiling}
          \resumeItemListEnd
      \resumeProjectHeading
          {\textbf{Useful Tool Hub} $|$ \emph{JavaScript, Python, Jest, Vercel} $|$ \href{https://useful-tool-hub.vercel.app}{Live}}{2026}
          \resumeItemListStart
            \resumeItem{Made a dozen web tools (file conversion, OCR, transcription, PDF signing) that run in the browser, with a Python backend only for downloaders; Jest runs in CI and scripts check the deployed site}
          \resumeItemListEnd
      \resumeProjectHeading
          {\textbf{Us (browser story game)} $|$ \emph{Phaser 4, TypeScript, Playwright} $|$ \href{https://martinw500.github.io/larissa-game/}{Live}}{2026}
          \resumeItemListStart
            \resumeItem{Created 15+ scenes in Phaser 4 with saves, fast travel and dialogue rewind; Playwright scripts play through routes and check lighting and collision}
          \resumeItemListEnd
    \resumeSubHeadingListEnd

\section{Skills \& Awards}
 \begin{itemize}[leftmargin=0.15in, label={}]
    \small{\item{
    \textbf{Languages}{: C, Python, Go, TypeScript, JavaScript, Java, SystemVerilog} \\
    \textbf{Embedded}{: STM32 HAL, FreeRTOS, CAN bus, SPI, ADC, CMake, oscilloscope, PCAN, J-Link, ST-Link} \\
    \textbf{Frameworks \& Tools}{: Linux, Git, React, Vite, Docker, Grafana, Playwright, Jest} \\
    \textbf{Awards}{: CSMC Top 0.5\% nationally (71st) $|$ AIME Qualifier 2x $|$ Euclid Top 2\% (85/100)}
    }}
 \end{itemize}

\end{document}
